% Page 41 — Alban au temps de François 1er

\begin{center}
  {\Large\itshape Alban au temps de François 1\textsuperscript{er} (XVI\textsuperscript{ème} siècle)}
\end{center}

\vspace{1em}

\lettrine{U}{ne} période paisible : cette première moitié du XVI\textsuperscript{ème} siècle. Pour notre coin de montagne, les archives nous livrent le nom de quelques personnalités et font état d’une réorganisation tant au point de vue politique que religieux de ce terroir.

\medskip

Parmi les ecclésiastiques de cette époque, on trouve comme Prieur de N.D. d’Alban en 1505 un Raymond Guittard ; un peu plus tard, en 1508, c’est un Bertrand de Solages. (Les de Solages ont encore aujourd’hui dans leur famille plusieurs prêtres : Monseigneur Bruno de Solages, ancien recteur de l’Institut Catholique de Toulouse, François de Solages, curé de St Jean de Tartage à Gaillac, Hugues de Solages, curé de Cuq et de Serviès, etc…). Un ancien curé de Courris est signalé comme prieur de N.D. d’Alban en 1551 ; il s’agit de Jean de Moret. Ces de Moret illustrent particulièrement notre histoire locale. Raymond Gausseran est bien nommé en 1545 comme prêtre à N.D. d’Alban, mais en était-il curé ? On ne sait.

\medskip

\lettrine{O}{n} désigne communément comme fondateur de la « chapellenie de St Laurent » en 1516, un « Louis de Moret ». Il semble qu’à cela il y ait une petite difficulté, puisque déjà, le 14 juin 1457, on trouve un Jean Cathala au « Mas de la Cappélanyé » (actuellement la ferme de Capel-le-Bas, sous la nouvelle gendarmerie). On peut peut-être penser que Louis de Moret n’a fait qu’améliorer une fondation existante en la dotant de biens plus importants. A ce moment-là (en 1516), la chapellenie joint d’un immeuble dans Alban même et aux alentours du « mas » d’un pré dit « de la Cappélanyé ». Ce dernier confronte le hameau de la Drèch alors appelé « Calm », la route d’Alban à St Sernin et le rec (ruisseau) de la Cappélanyé. Les revenus de ces biens servaient à payer les honoraires des messes acquittées pour la famille des de Moret de Nougaret dont l’un d’eux était seigneur de la baronnie d’Alban.
